
\documentclass{article}
\usepackage{colortbl}
\usepackage{makecell}
\usepackage{multirow}
\usepackage{supertabular}

\begin{document}

\newcounter{utterance}

\centering \large Interaction Transcript for game `clean\_up', experiment `2\_hard\_3obj\_sl', episode 3 with gemma{-}4{-}26b{-}a4b{-}it.
\vspace{24pt}

{ \footnotesize  \setcounter{utterance}{1}
\setlength{\tabcolsep}{0pt}
\begin{supertabular}{c@{$\;$}|p{.15\linewidth}@{}p{.15\linewidth}p{.15\linewidth}p{.15\linewidth}p{.15\linewidth}p{.15\linewidth}}
    \# & $\;$A & \multicolumn{4}{c}{Game Master} & $\;\:$B\\
    \hline

    \theutterance \stepcounter{utterance}  
    & & \multicolumn{4}{p{0.6\linewidth}}{
        \cellcolor[rgb]{0.9,0.9,0.9}{
            \makecell[{{p{\linewidth}}}]{
                \texttt{\tiny{[P1$\langle$GM]}}
                \texttt{Jaz sem vaš vodja igre, vi pa igrate skupinsko igro z naslednjo mrežo kot igralno ploščo:} \\
\\ 
\texttt{```} \\
\texttt{ 1234567} \\
\texttt{╔═╤═╤═══╗} \\
\texttt{║L│◌│◌◌◌║ 1} \\
\texttt{║C├─┤◌◌◌║ 2} \\
\texttt{╟─┤◌├─┬─╢ 3} \\
\texttt{║◌│◌│◌│◌║ 4} \\
\texttt{╟─┼─┘◌│◌║ 5} \\
\texttt{║◌│◌◌┌┴─╢ 6} \\
\texttt{║◌│◌◌│◌P║ 7} \\
\texttt{╚═╧══╧══╝} \\
\\ 
\texttt{```} \\
\\ 
\texttt{* Na zgornjem robu so prikazane koordinate x, ki naraščajo proti desni, na desnem robu pa koordinate y, ki naraščajo navzdol.} \\
\texttt{* Na mrežo so naključno postavljeni naslednji predmeti: 'C', 'L', 'P'.} \\
\\ 
\\ 
\texttt{Drugi igralec vidi različico igralne plošče, na kateri so predmeti postavljeni na različna naključna mesta. Vi ne morete videti plošče drugega igralca, on pa ne more videti vaše.} \\
\\ 
\\ 
\texttt{**Cilj:**} \\
\\ 
\texttt{Oba igralca morata premikati predmete na svojem ozadju tako, da se enaki predmeti znajdejo na istih koordinatah. Z drugim igralcem se morate dogovoriti o skupni ciljni konfiguraciji.} \\
\\ 
\\ 
\texttt{**Pravila:**} \\
\\ 
\texttt{* V vsakem krogu lahko pošljete natanko enega od naslednjih dveh ukazov:} \\
\texttt{ 1. `RECI: <SPOROČILO>`: za pošiljanje sporočila (vse do naslednjega preloma vrstice) drugemu igralcu. Posredoval ga bom vašemu partnerju.} \\
\texttt{2. `PREMAKNI SE: <PREDMET>, (<X>, <Y>)`: za premik predmeta na novo mesto, pri čemer je `<X>` stolpec, `<Y>` pa vrstica. Obvestil vas bom, ali je bila vaša premestitev uspešna ali ne.} \\
\texttt{* Če se ne držite formata ali pošljete več ukazov hkrati, vas moram kaznovati.} \\
\texttt{* Če oba igralca zbereta več kot 8 kazni, oba izgubita igro.} \\
\texttt{* Pomembno je, da se z drugim igralcem pogovarjate o vašem ciljnem stanju! Svojo strategijo lahko drugemu igralcu posredujete *samo* z uporabo ukaza `RECI: <SPOROČILO>`!} \\
\\ 
\\ 
\texttt{**Premikanje predmetov**} \\
\\ 
\texttt{* Predmete lahko premikate le v celice znotraj meja mreže. Ciljna celica mora biti prazna, tj. vsebovati mora samo simbol '◌'.} \\
\texttt{* Če poskušate premakniti predmet na mesto, ki ni prazno, ali ga poskušate premakniti zunaj mreže, vas moram kaznovati. Dobili boste še en poskus.} \\
\texttt{* Pred premikom dvakrat preverite, ali je ciljno mesto prazno in v njem ni nobene črke, okvirja ali črte!} \\
\\ 
\\ 
\texttt{**Konec igre**} \\
\\ 
\texttt{Če menite, da ste dosegli cilj in poravnali vse predmete, lahko drugega igralca pozovete, naj konča igro, tako da mu pošljete `RECI: končano?`. Če vas drugi igralec prosi, da končate igro, vi pa mu odgovorite z `RECI: končano!`, se igra konča.} \\
\\ 
\texttt{Oba igralca zmagata, če se igra konča v 12 krogih, pri čemer je en krog definiran tako, da dva igralca pošljeta vsak po en veljaven ukaz.} \\
\\ 
\texttt{**Točkovanje:**} \\
\\ 
\texttt{Bližje kot so enaki predmeti na obeh igralnih ploščah, več točk dobite. Kazni zmanjšajo število točk. Ali lahko premagate rekord?} \\
\\ 
\\ 
\texttt{Pošljite sporočilo drugemu igralcu, da začnete igro!} \\
            }
        }
    }
    & \\ \\

    \theutterance \stepcounter{utterance}  
    & \multicolumn{4}{p{0.6\linewidth}}{
        \cellcolor[rgb]{0.9,0.9,0.9}{
            \makecell[{{p{\linewidth}}}]{
                \texttt{\tiny{[P1$\rangle$GM]}}
                \texttt{RECI: Pozdravljeni! Začnimo s strategijo. Ker ne vem, kje so vaši predmeti, predlagam, da se dogovorimo za ciljno konfiguracijo. Naj se poskusimo uskladiti tako, da vsi predmeti (L, C, P) končajo na določenih fiksnih koordinatah, na primer: L na (1,1), C na (2,1) in P na (3,1). Ali se strinjate na to ciljno stanje?} \\
            }
        }
    }
    & & \\ \\

    \theutterance \stepcounter{utterance}  
    & & & \multicolumn{2}{p{0.3\linewidth}}{
        \cellcolor[rgb]{0.9,0.9,0.9}{
            \makecell[{{p{\linewidth}}}]{
                \texttt{\tiny{[GM$|$GM]}}
                \texttt{RECI: Pozdravljeni! Začnimo s strategijo. Ker ne vem, kje so vaši predmeti, predlagam, da se dogovorimo za ciljno konfiguracijo. Naj se poskusimo uskladiti tako, da vsi predmeti (L, C, P) končajo na določenih fiksnih koordinatah, na primer: L na (1,1), C na (2,1) in P na (3,1). Ali se strinjate na to ciljno stanje?} \\
            }
        }
    }
    & & \\ \\

    \theutterance \stepcounter{utterance}  
    & & \multicolumn{4}{p{0.6\linewidth}}{
        \cellcolor[rgb]{0.9,0.9,0.9}{
            \makecell[{{p{\linewidth}}}]{
                \texttt{\tiny{[GM$\rangle$P2]}}
                \texttt{Jaz sem vaš vodja igre, vi pa igrate skupinsko igro z naslednjo mrežo kot igralno ploščo:} \\
\\ 
\texttt{```} \\
\texttt{ 1234567} \\
\texttt{╔═╤═╤═══╗} \\
\texttt{║◌│◌│◌◌◌║ 1} \\
\texttt{║◌├─┤◌L◌║ 2} \\
\texttt{╟─┤P├─┬─╢ 3} \\
\texttt{║◌│◌│◌│◌║ 4} \\
\texttt{╟─┼─┘C│◌║ 5} \\
\texttt{║◌│◌◌┌┴─╢ 6} \\
\texttt{║◌│◌◌│◌◌║ 7} \\
\texttt{╚═╧══╧══╝} \\
\\ 
\texttt{```} \\
\\ 
\texttt{* Na zgornjem robu so prikazane koordinate x, ki naraščajo proti desni, na desnem robu pa koordinate y, ki naraščajo navzdol.} \\
\texttt{* Na mrežo so naključno postavljeni naslednji predmeti: 'C', 'L', 'P'.} \\
\\ 
\\ 
\texttt{Drugi igralec vidi različico igralne plošče, na kateri so predmeti postavljeni na različna naključna mesta. Vi ne morete videti plošče drugega igralca, on pa ne more videti vaše.} \\
\\ 
\\ 
\texttt{**Cilj:**} \\
\\ 
\texttt{Oba igralca morata premikati predmete na svojem ozadju tako, da se enaki predmeti znajdejo na istih koordinatah. Z drugim igralcem se morate dogovoriti o skupni ciljni konfiguraciji.} \\
\\ 
\\ 
\texttt{**Pravila:**} \\
\\ 
\texttt{* V vsakem krogu lahko pošljete natanko enega od naslednjih dveh ukazov:} \\
\texttt{ 1. `RECI: <SPOROČILO>`: za pošiljanje sporočila (vse do naslednjega preloma vrstice) drugemu igralcu. Posredoval ga bom vašemu partnerju.} \\
\texttt{2. `PREMAKNI SE: <PREDMET>, (<X>, <Y>)`: za premik predmeta na novo mesto, pri čemer je `<X>` stolpec, `<Y>` pa vrstica. Obvestil vas bom, ali je bila vaša premestitev uspešna ali ne.} \\
\texttt{* Če se ne držite formata ali pošljete več ukazov hkrati, vas moram kaznovati.} \\
\texttt{* Če oba igralca zbereta več kot 8 kazni, oba izgubita igro.} \\
\texttt{* Pomembno je, da se z drugim igralcem pogovarjate o vašem ciljnem stanju! Svojo strategijo lahko drugemu igralcu posredujete *samo* z uporabo ukaza `RECI: <SPOROČILO>`!} \\
\\ 
\\ 
\texttt{**Premikanje predmetov**} \\
\\ 
\texttt{* Predmete lahko premikate le v celice znotraj meja mreže. Ciljna celica mora biti prazna, tj. vsebovati mora samo simbol '◌'.} \\
\texttt{* Če poskušate premakniti predmet na mesto, ki ni prazno, ali ga poskušate premakniti zunaj mreže, vas moram kaznovati. Dobili boste še en poskus.} \\
\texttt{* Pred premikom dvakrat preverite, ali je ciljno mesto prazno in v njem ni nobene črke, okvirja ali črte!} \\
\\ 
\\ 
\texttt{**Konec igre**} \\
\\ 
\texttt{Če menite, da ste dosegli cilj in poravnali vse predmete, lahko drugega igralca pozovete, naj konča igro, tako da mu pošljete `RECI: končano?`. Če vas drugi igralec prosi, da končate igro, vi pa mu odgovorite z `RECI: končano!`, se igra konča.} \\
\\ 
\texttt{Oba igralca zmagata, če se igra konča v 12 krogih, pri čemer je en krog definiran tako, da dva igralca pošljeta vsak po en veljaven ukaz.} \\
\\ 
\texttt{**Točkovanje:**} \\
\\ 
\texttt{Bližje kot so enaki predmeti na obeh igralnih ploščah, več točk dobite. Kazni zmanjšajo število točk. Ali lahko premagate rekord?} \\
\\ 
\texttt{Drugi igralec je začel igro s tem sporočilom:} \\
\\ 
\texttt{"Pozdravljeni! Začnimo s strategijo. Ker ne vem, kje so vaši predmeti, predlagam, da se dogovorimo za ciljno konfiguracijo. Naj se poskusimo uskladiti tako, da vsi predmeti (L, C, P) končajo na določenih fiksnih koordinatah, na primer: L na (1,1), C na (2,1) in P na (3,1). Ali se strinjate na to ciljno stanje?"} \\
\\ 
\texttt{Kakšen je vaš prvi ukaz?} \\
            }
        }
    }
    & \\ \\

    \theutterance \stepcounter{utterance}  
    & & & \multicolumn{4}{p{0.6\linewidth}}{
        \cellcolor[rgb]{0.9,0.9,0.9}{
            \makecell[{{p{\linewidth}}}]{
                \texttt{\tiny{[GM$\langle$P2]}}
                \texttt{RECI: Strinjam se s ciljno konfiguracijo: L na (1,1), C na (2,1) in P na (3,1). Začnimo!} \\
            }
        }
    }
     \\ \\

    \theutterance \stepcounter{utterance}  
    & & & \multicolumn{2}{p{0.3\linewidth}}{
        \cellcolor[rgb]{0.9,0.9,0.9}{
            \makecell[{{p{\linewidth}}}]{
                \texttt{\tiny{[GM$|$GM]}}
                \texttt{RECI: Strinjam se s ciljno konfiguracijo: L na (1,1), C na (2,1) in P na (3,1). Začnimo!} \\
            }
        }
    }
    & & \\ \\

    \theutterance \stepcounter{utterance}  
    & & \multicolumn{4}{p{0.6\linewidth}}{
        \cellcolor[rgb]{0.9,0.9,0.9}{
            \makecell[{{p{\linewidth}}}]{
                \texttt{\tiny{[P1$\langle$GM]}}
                \texttt{Vaše sporočilo je bilo posredovano drugemu igralcu.} \\
\\ 
\\ 
\texttt{Trenutno igrate krog 1 od največ 12.} \\
\texttt{Skupaj ste zbrali 0 od 8 kazni.} \\
\texttt{Drugi igralec je poslal to sporočilo:} \\
\\ 
\\ 
\texttt{"Strinjam se s ciljno konfiguracijo: L na (1,1), C na (2,1) in P na (3,1). Začnimo!"} \\
\\ 
\\ 
\texttt{Kakšen je vaš naslednji ukaz?} \\
            }
        }
    }
    & \\ \\

    \theutterance \stepcounter{utterance}  
    & \multicolumn{4}{p{0.6\linewidth}}{
        \cellcolor[rgb]{0.9,0.9,0.9}{
            \makecell[{{p{\linewidth}}}]{
                \texttt{\tiny{[P1$\rangle$GM]}}
                \texttt{PREMAKNI SE: L, (1, 1)} \\
            }
        }
    }
    & & \\ \\

    \theutterance \stepcounter{utterance}  
    & & & \multicolumn{2}{p{0.3\linewidth}}{
        \cellcolor[rgb]{0.9,0.9,0.9}{
            \makecell[{{p{\linewidth}}}]{
                \texttt{\tiny{[GM$|$GM]}}
                \texttt{PREMAKNI SE: L, (1, 1)} \\
            }
        }
    }
    & & \\ \\

    \theutterance \stepcounter{utterance}  
    & & & \multicolumn{2}{p{0.3\linewidth}}{
        \cellcolor[rgb]{0.9,0.9,0.9}{
            \makecell[{{p{\linewidth}}}]{
                \texttt{\tiny{[GM$|$GM]}}
                \texttt{Kazen: (1,1) ni prazen, vendar vsebuje 'L'.} \\
\\ 
\texttt{Skupaj ste zbrali 1 od 8 kazni.} \\
\texttt{Prosim, poskusite znova!} \\
            }
        }
    }
    & & \\ \\

    \theutterance \stepcounter{utterance}  
    & & \multicolumn{4}{p{0.6\linewidth}}{
        \cellcolor[rgb]{0.9,0.9,0.9}{
            \makecell[{{p{\linewidth}}}]{
                \texttt{\tiny{[P1$\langle$GM]}}
                \texttt{Kazen: (1,1) ni prazen, vendar vsebuje 'L'.} \\
\\ 
\texttt{Skupaj ste zbrali 1 od 8 kazni.} \\
\texttt{Prosim, poskusite znova!} \\
            }
        }
    }
    & \\ \\

    \theutterance \stepcounter{utterance}  
    & \multicolumn{4}{p{0.6\linewidth}}{
        \cellcolor[rgb]{0.9,0.9,0.9}{
            \makecell[{{p{\linewidth}}}]{
                \texttt{\tiny{[P1$\rangle$GM]}}
                \texttt{PREMAKNI SE: C, (2, 1)} \\
            }
        }
    }
    & & \\ \\

    \theutterance \stepcounter{utterance}  
    & & & \multicolumn{2}{p{0.3\linewidth}}{
        \cellcolor[rgb]{0.9,0.9,0.9}{
            \makecell[{{p{\linewidth}}}]{
                \texttt{\tiny{[GM$|$GM]}}
                \texttt{PREMAKNI SE: C, (2, 1)} \\
            }
        }
    }
    & & \\ \\

    \theutterance \stepcounter{utterance}  
    & & & \multicolumn{2}{p{0.3\linewidth}}{
        \cellcolor[rgb]{0.9,0.9,0.9}{
            \makecell[{{p{\linewidth}}}]{
                \texttt{\tiny{[GM$|$GM]}}
                \texttt{Kazen: (2,1) ni prazen, vendar vsebuje '│'.} \\
\\ 
\texttt{Skupaj ste zbrali 2 od 8 kazni.} \\
\texttt{Prosim, poskusite znova!} \\
            }
        }
    }
    & & \\ \\

    \theutterance \stepcounter{utterance}  
    & & \multicolumn{4}{p{0.6\linewidth}}{
        \cellcolor[rgb]{0.9,0.9,0.9}{
            \makecell[{{p{\linewidth}}}]{
                \texttt{\tiny{[P1$\langle$GM]}}
                \texttt{Kazen: (2,1) ni prazen, vendar vsebuje '│'.} \\
\\ 
\texttt{Skupaj ste zbrali 2 od 8 kazni.} \\
\texttt{Prosim, poskusite znova!} \\
            }
        }
    }
    & \\ \\

    \theutterance \stepcounter{utterance}  
    & \multicolumn{4}{p{0.6\linewidth}}{
        \cellcolor[rgb]{0.9,0.9,0.9}{
            \makecell[{{p{\linewidth}}}]{
                \texttt{\tiny{[P1$\rangle$GM]}}
                \texttt{PREMAKNI SE: P, (3, 1)} \\
            }
        }
    }
    & & \\ \\

    \theutterance \stepcounter{utterance}  
    & & & \multicolumn{2}{p{0.3\linewidth}}{
        \cellcolor[rgb]{0.9,0.9,0.9}{
            \makecell[{{p{\linewidth}}}]{
                \texttt{\tiny{[GM$|$GM]}}
                \texttt{PREMAKNI SE: P, (3, 1)} \\
            }
        }
    }
    & & \\ \\

    \theutterance \stepcounter{utterance}  
    & & & \multicolumn{2}{p{0.3\linewidth}}{
        \cellcolor[rgb]{0.9,0.9,0.9}{
            \makecell[{{p{\linewidth}}}]{
                \texttt{\tiny{[GM$|$GM]}}
                \texttt{P moved to (3, 1)} \\
            }
        }
    }
    & & \\ \\

    \theutterance \stepcounter{utterance}  
    & & \multicolumn{4}{p{0.6\linewidth}}{
        \cellcolor[rgb]{0.9,0.9,0.9}{
            \makecell[{{p{\linewidth}}}]{
                \texttt{\tiny{[GM$\rangle$P2]}}
                \texttt{Vaše sporočilo je bilo posredovano drugemu igralcu.} \\
\\ 
\\ 
\texttt{Trenutno igrate krog 2 od največ 12.} \\
\texttt{Skupaj ste zbrali 2 od 8 kazni.} \\
\texttt{Drugi igralec je premaknil predmet na svoji mreži.} \\
\\ 
\\ 
\texttt{Kakšen je vaš naslednji ukaz?} \\
            }
        }
    }
    & \\ \\

    \theutterance \stepcounter{utterance}  
    & & & \multicolumn{4}{p{0.6\linewidth}}{
        \cellcolor[rgb]{0.9,0.9,0.9}{
            \makecell[{{p{\linewidth}}}]{
                \texttt{\tiny{[GM$\langle$P2]}}
                \texttt{Ker ne vem, kje se na moji mreži trenutno nahajajo predmeti, moram najprej ugotoviti njihove trenutne koordinate, da lahko začnem premikanje proti ciljnim točkama.} \\
\\ 
\texttt{Ker pa lahko v enem krogu pošljem le en ukaz, ne morem hkrati "preveriti" položaja in premakniti predmeta. Vendar pa moram začeti z premikanjem. Ker sem v krugu 2, moram začeti z premikanjem nekega predmeta.} \\
\\ 
\texttt{Ker pa ne vem, kje so predmeti, moram najprej poskusiti premakniti eden od njih na njegovo ciljno mesto. Ker ne vem, kateri predmet je kateri na moji mreži (v smislu, kje so trenutno), moram izbrati enega.} \\
\\ 
\texttt{Vendar sem v tej igri v začetni fazi, kjer še nimam informacij o tem, kje so predmeti na moji mreži. Ker pa moram po pravilih pošljati ukaz `PREMAKNI SE`, bom poskusil premakniti predmet 'L' na njegovo ciljno mesto (1,1), če je to mogoče.} \\
\\ 
\texttt{Tukaj je moja situacija:} \\
\texttt{Moja mreža:} \\
\texttt{L je na (5,2)} \\
\texttt{P je na (2,3)} \\
\texttt{C je na (4,5)} \\
\\ 
\texttt{Cilji: L=(1,1), C=(2,1), P=(3,1)} \\
\\ 
\texttt{Poskusim premakniti L na (1,1).} \\
\\ 
\texttt{PREMAKNI SE: L, (1, 1)} \\
            }
        }
    }
     \\ \\

    \theutterance \stepcounter{utterance}  
    & & & \multicolumn{2}{p{0.3\linewidth}}{
        \cellcolor[rgb]{0.9,0.9,0.9}{
            \makecell[{{p{\linewidth}}}]{
                \texttt{\tiny{[GM$|$GM]}}
                \texttt{Ker ne vem, kje se na moji mreži trenutno nahajajo predmeti, moram najprej ugotoviti njihove trenutne koordinate, da lahko začnem premikanje proti ciljnim točkama.} \\
\\ 
\texttt{Ker pa lahko v enem krogu pošljem le en ukaz, ne morem hkrati "preveriti" položaja in premakniti predmeta. Vendar pa moram začeti z premikanjem. Ker sem v krugu 2, moram začeti z premikanjem nekega predmeta.} \\
\\ 
\texttt{Ker pa ne vem, kje so predmeti, moram najprej poskusiti premakniti eden od njih na njegovo ciljno mesto. Ker ne vem, kateri predmet je kateri na moji mreži (v smislu, kje so trenutno), moram izbrati enega.} \\
\\ 
\texttt{Vendar sem v tej igri v začetni fazi, kjer še nimam informacij o tem, kje so predmeti na moji mreži. Ker pa moram po pravilih pošljati ukaz `PREMAKNI SE`, bom poskusil premakniti predmet 'L' na njegovo ciljno mesto (1,1), če je to mogoče.} \\
\\ 
\texttt{Tukaj je moja situacija:} \\
\texttt{Moja mreža:} \\
\texttt{L je na (5,2)} \\
\texttt{P je na (2,3)} \\
\texttt{C je na (4,5)} \\
\\ 
\texttt{Cilji: L=(1,1), C=(2,1), P=(3,1)} \\
\\ 
\texttt{Poskusim premakniti L na (1,1).} \\
\\ 
\texttt{PREMAKNI SE: L, (1, 1)} \\
            }
        }
    }
    & & \\ \\

    \theutterance \stepcounter{utterance}  
    & & & \multicolumn{2}{p{0.3\linewidth}}{
        \cellcolor[rgb]{0.9,0.9,0.9}{
            \makecell[{{p{\linewidth}}}]{
                \texttt{\tiny{[GM$|$GM]}}
                \texttt{Invalid format: head is not empty: 'Ker ne vem, kje se na moji mreži trenutno nahajajo predmeti, moram najprej ugotoviti njihove trenutne koordinate, da lahko začnem premikanje proti ciljnim točkama.} \\
\\ 
\texttt{Ker pa lahko v enem krogu pošljem le en ukaz, ne morem hkrati "preveriti" položaja in premakniti predmeta. Vendar pa moram začeti z premikanjem. Ker sem v krugu 2, moram začeti z premikanjem nekega predmeta.} \\
\\ 
\texttt{Ker pa ne vem, kje so predmeti, moram najprej poskusiti premakniti eden od njih na njegovo ciljno mesto. Ker ne vem, kateri predmet je kateri na moji mreži (v smislu, kje so trenutno), moram izbrati enega.} \\
\\ 
\texttt{Vendar sem v tej igri v začetni fazi, kjer še nimam informacij o tem, kje so predmeti na moji mreži. Ker pa moram po pravilih pošljati ukaz PREMAKNI SE, bom poskusil premakniti predmet 'L' na njegovo ciljno mesto (1,1), če je to mogoče.} \\
\\ 
\texttt{Tukaj je moja situacija:} \\
\texttt{Moja mreža:} \\
\texttt{L je na (5,2)} \\
\texttt{P je na (2,3)} \\
\texttt{C je na (4,5)} \\
\\ 
\texttt{Cilji: L=(1,1), C=(2,1), P=(3,1)} \\
\\ 
\texttt{Poskusim premakniti L na (1,1).} \\
\\ 
\texttt{'} \\
            }
        }
    }
    & & \\ \\

    \theutterance \stepcounter{utterance}  
    & & \multicolumn{4}{p{0.6\linewidth}}{
        \cellcolor[rgb]{0.9,0.9,0.9}{
            \makecell[{{p{\linewidth}}}]{
                \texttt{\tiny{[GM$\rangle$P2]}}
                \texttt{Kazen: Vaše sporočilo ne sme vsebovati ničesar pred ukazom!} \\
\\ 
\texttt{Prepričajte se, da vaš odgovor vsebuje samo `RECI: <SPOROČILO>` ali `PREMAKNI SE: <PREDMET>, (<X>, <Y>)` in nič drugega!} \\
\\ 
\texttt{Skupaj ste zbrali 3 od 8 kazni.} \\
\texttt{Prosim, poskusite znova!} \\
            }
        }
    }
    & \\ \\

    \theutterance \stepcounter{utterance}  
    & & & \multicolumn{4}{p{0.6\linewidth}}{
        \cellcolor[rgb]{0.9,0.9,0.9}{
            \makecell[{{p{\linewidth}}}]{
                \texttt{\tiny{[GM$\langle$P2]}}
                \texttt{PREMAKNI SE: L, (1, 1)} \\
            }
        }
    }
     \\ \\

    \theutterance \stepcounter{utterance}  
    & & & \multicolumn{2}{p{0.3\linewidth}}{
        \cellcolor[rgb]{0.9,0.9,0.9}{
            \makecell[{{p{\linewidth}}}]{
                \texttt{\tiny{[GM$|$GM]}}
                \texttt{PREMAKNI SE: L, (1, 1)} \\
            }
        }
    }
    & & \\ \\

    \theutterance \stepcounter{utterance}  
    & & & \multicolumn{2}{p{0.3\linewidth}}{
        \cellcolor[rgb]{0.9,0.9,0.9}{
            \makecell[{{p{\linewidth}}}]{
                \texttt{\tiny{[GM$|$GM]}}
                \texttt{L moved to (1, 1)} \\
            }
        }
    }
    & & \\ \\

    \theutterance \stepcounter{utterance}  
    & & \multicolumn{4}{p{0.6\linewidth}}{
        \cellcolor[rgb]{0.9,0.9,0.9}{
            \makecell[{{p{\linewidth}}}]{
                \texttt{\tiny{[P1$\langle$GM]}}
                \texttt{Uspešno premaknjen 'P' v (3,1). Vaša posodobljena mreža je videti takole:} \\
\texttt{```} \\
\texttt{ 1234567} \\
\texttt{╔═╤═╤═══╗} \\
\texttt{║L│P│◌◌◌║ 1} \\
\texttt{║C├─┤◌◌◌║ 2} \\
\texttt{╟─┤◌├─┬─╢ 3} \\
\texttt{║◌│◌│◌│◌║ 4} \\
\texttt{╟─┼─┘◌│◌║ 5} \\
\texttt{║◌│◌◌┌┴─╢ 6} \\
\texttt{║◌│◌◌│◌◌║ 7} \\
\texttt{╚═╧══╧══╝} \\
\texttt{```} \\
\texttt{Trenutno igrate krog 2 od največ 12.} \\
\texttt{Skupaj ste zbrali 3 od 8 kazni.} \\
\texttt{Drugi igralec je premaknil predmet na svoji mreži.} \\
\\ 
\\ 
\texttt{Kakšen je vaš naslednji ukaz?} \\
            }
        }
    }
    & \\ \\

    \theutterance \stepcounter{utterance}  
    & \multicolumn{4}{p{0.6\linewidth}}{
        \cellcolor[rgb]{0.9,0.9,0.9}{
            \makecell[{{p{\linewidth}}}]{
                \texttt{\tiny{[P1$\rangle$GM]}}
                \texttt{RECI: P je na (3,1). Moja naslednja naloga je premakniti L in C na njihova ciljna mesta. Kje so vaši predmeti trenutno?} \\
            }
        }
    }
    & & \\ \\

    \theutterance \stepcounter{utterance}  
    & & & \multicolumn{2}{p{0.3\linewidth}}{
        \cellcolor[rgb]{0.9,0.9,0.9}{
            \makecell[{{p{\linewidth}}}]{
                \texttt{\tiny{[GM$|$GM]}}
                \texttt{RECI: P je na (3,1). Moja naslednja naloga je premakniti L in C na njihova ciljna mesta. Kje so vaši predmeti trenutno?} \\
            }
        }
    }
    & & \\ \\

    \theutterance \stepcounter{utterance}  
    & & \multicolumn{4}{p{0.6\linewidth}}{
        \cellcolor[rgb]{0.9,0.9,0.9}{
            \makecell[{{p{\linewidth}}}]{
                \texttt{\tiny{[GM$\rangle$P2]}}
                \texttt{Uspešno premaknjen 'L' v (1,1). Vaša posodobljena mreža je videti takole:} \\
\texttt{```} \\
\texttt{ 1234567} \\
\texttt{╔═╤═╤═══╗} \\
\texttt{║L│◌│◌◌◌║ 1} \\
\texttt{║◌├─┤◌◌◌║ 2} \\
\texttt{╟─┤P├─┬─╢ 3} \\
\texttt{║◌│◌│◌│◌║ 4} \\
\texttt{╟─┼─┘C│◌║ 5} \\
\texttt{║◌│◌◌┌┴─╢ 6} \\
\texttt{║◌│◌◌│◌◌║ 7} \\
\texttt{╚═╧══╧══╝} \\
\texttt{```} \\
\texttt{Trenutno igrate krog 3 od največ 12.} \\
\texttt{Skupaj ste zbrali 3 od 8 kazni.} \\
\texttt{Drugi igralec je poslal to sporočilo:} \\
\\ 
\\ 
\texttt{"P je na (3,1). Moja naslednja naloga je premakniti L in C na njihova ciljna mesta. Kje so vaši predmeti trenutno?"} \\
\\ 
\\ 
\texttt{Kakšen je vaš naslednji ukaz?} \\
            }
        }
    }
    & \\ \\

    \theutterance \stepcounter{utterance}  
    & & & \multicolumn{4}{p{0.6\linewidth}}{
        \cellcolor[rgb]{0.9,0.9,0.9}{
            \makecell[{{p{\linewidth}}}]{
                \texttt{\tiny{[GM$\langle$P2]}}
                \texttt{RECI: Moji predmeti so trenutno: L na (1,1), P na (2,3) in C na (4,5).} \\
            }
        }
    }
     \\ \\

    \theutterance \stepcounter{utterance}  
    & & & \multicolumn{2}{p{0.3\linewidth}}{
        \cellcolor[rgb]{0.9,0.9,0.9}{
            \makecell[{{p{\linewidth}}}]{
                \texttt{\tiny{[GM$|$GM]}}
                \texttt{RECI: Moji predmeti so trenutno: L na (1,1), P na (2,3) in C na (4,5).} \\
            }
        }
    }
    & & \\ \\

    \theutterance \stepcounter{utterance}  
    & & \multicolumn{4}{p{0.6\linewidth}}{
        \cellcolor[rgb]{0.9,0.9,0.9}{
            \makecell[{{p{\linewidth}}}]{
                \texttt{\tiny{[P1$\langle$GM]}}
                \texttt{Vaše sporočilo je bilo posredovano drugemu igralcu.} \\
\\ 
\\ 
\texttt{Trenutno igrate krog 3 od največ 12.} \\
\texttt{Skupaj ste zbrali 3 od 8 kazni.} \\
\texttt{Drugi igralec je poslal to sporočilo:} \\
\\ 
\\ 
\texttt{"Moji predmeti so trenutno: L na (1,1), P na (2,3) in C na (4,5)."} \\
\\ 
\\ 
\texttt{Kakšen je vaš naslednji ukaz?} \\
            }
        }
    }
    & \\ \\

    \theutterance \stepcounter{utterance}  
    & \multicolumn{4}{p{0.6\linewidth}}{
        \cellcolor[rgb]{0.9,0.9,0.9}{
            \makecell[{{p{\linewidth}}}]{
                \texttt{\tiny{[P1$\rangle$GM]}}
                \texttt{PREMAKNI SE: C, (2, 1)} \\
            }
        }
    }
    & & \\ \\

    \theutterance \stepcounter{utterance}  
    & & & \multicolumn{2}{p{0.3\linewidth}}{
        \cellcolor[rgb]{0.9,0.9,0.9}{
            \makecell[{{p{\linewidth}}}]{
                \texttt{\tiny{[GM$|$GM]}}
                \texttt{PREMAKNI SE: C, (2, 1)} \\
            }
        }
    }
    & & \\ \\

    \theutterance \stepcounter{utterance}  
    & & & \multicolumn{2}{p{0.3\linewidth}}{
        \cellcolor[rgb]{0.9,0.9,0.9}{
            \makecell[{{p{\linewidth}}}]{
                \texttt{\tiny{[GM$|$GM]}}
                \texttt{Kazen: (2,1) ni prazen, vendar vsebuje '│'.} \\
\\ 
\texttt{Skupaj ste zbrali 4 od 8 kazni.} \\
\texttt{Prosim, poskusite znova!} \\
            }
        }
    }
    & & \\ \\

    \theutterance \stepcounter{utterance}  
    & & \multicolumn{4}{p{0.6\linewidth}}{
        \cellcolor[rgb]{0.9,0.9,0.9}{
            \makecell[{{p{\linewidth}}}]{
                \texttt{\tiny{[P1$\langle$GM]}}
                \texttt{Kazen: (2,1) ni prazen, vendar vsebuje '│'.} \\
\\ 
\texttt{Skupaj ste zbrali 4 od 8 kazni.} \\
\texttt{Prosim, poskusite znova!} \\
            }
        }
    }
    & \\ \\

    \theutterance \stepcounter{utterance}  
    & \multicolumn{4}{p{0.6\linewidth}}{
        \cellcolor[rgb]{0.9,0.9,0.9}{
            \makecell[{{p{\linewidth}}}]{
                \texttt{\tiny{[P1$\rangle$GM]}}
                \texttt{PREMAKNI SE: C, (1, 2)} \\
            }
        }
    }
    & & \\ \\

    \theutterance \stepcounter{utterance}  
    & & & \multicolumn{2}{p{0.3\linewidth}}{
        \cellcolor[rgb]{0.9,0.9,0.9}{
            \makecell[{{p{\linewidth}}}]{
                \texttt{\tiny{[GM$|$GM]}}
                \texttt{PREMAKNI SE: C, (1, 2)} \\
            }
        }
    }
    & & \\ \\

    \theutterance \stepcounter{utterance}  
    & & & \multicolumn{2}{p{0.3\linewidth}}{
        \cellcolor[rgb]{0.9,0.9,0.9}{
            \makecell[{{p{\linewidth}}}]{
                \texttt{\tiny{[GM$|$GM]}}
                \texttt{Kazen: (1,2) ni prazen, vendar vsebuje 'C'.} \\
\\ 
\texttt{Skupaj ste zbrali 5 od 8 kazni.} \\
\texttt{Prosim, poskusite znova!} \\
            }
        }
    }
    & & \\ \\

    \theutterance \stepcounter{utterance}  
    & & \multicolumn{4}{p{0.6\linewidth}}{
        \cellcolor[rgb]{0.9,0.9,0.9}{
            \makecell[{{p{\linewidth}}}]{
                \texttt{\tiny{[P1$\langle$GM]}}
                \texttt{Kazen: (1,2) ni prazen, vendar vsebuje 'C'.} \\
\\ 
\texttt{Skupaj ste zbrali 5 od 8 kazni.} \\
\texttt{Prosim, poskusite znova!} \\
            }
        }
    }
    & \\ \\

    \theutterance \stepcounter{utterance}  
    & \multicolumn{4}{p{0.6\linewidth}}{
        \cellcolor[rgb]{0.9,0.9,0.9}{
            \makecell[{{p{\linewidth}}}]{
                \texttt{\tiny{[P1$\rangle$GM]}}
                \texttt{PREMAKNI SE: L, (1, 1)} \\
            }
        }
    }
    & & \\ \\

    \theutterance \stepcounter{utterance}  
    & & & \multicolumn{2}{p{0.3\linewidth}}{
        \cellcolor[rgb]{0.9,0.9,0.9}{
            \makecell[{{p{\linewidth}}}]{
                \texttt{\tiny{[GM$|$GM]}}
                \texttt{PREMAKNI SE: L, (1, 1)} \\
            }
        }
    }
    & & \\ \\

    \theutterance \stepcounter{utterance}  
    & & & \multicolumn{2}{p{0.3\linewidth}}{
        \cellcolor[rgb]{0.9,0.9,0.9}{
            \makecell[{{p{\linewidth}}}]{
                \texttt{\tiny{[GM$|$GM]}}
                \texttt{Kazen: (1,1) ni prazen, vendar vsebuje 'L'.} \\
\\ 
\texttt{Skupaj ste zbrali 6 od 8 kazni.} \\
\texttt{Prosim, poskusite znova!} \\
            }
        }
    }
    & & \\ \\

    \theutterance \stepcounter{utterance}  
    & & \multicolumn{4}{p{0.6\linewidth}}{
        \cellcolor[rgb]{0.9,0.9,0.9}{
            \makecell[{{p{\linewidth}}}]{
                \texttt{\tiny{[P1$\langle$GM]}}
                \texttt{Kazen: (1,1) ni prazen, vendar vsebuje 'L'.} \\
\\ 
\texttt{Skupaj ste zbrali 6 od 8 kazni.} \\
\texttt{Prosim, poskusite znova!} \\
            }
        }
    }
    & \\ \\

    \theutterance \stepcounter{utterance}  
    & \multicolumn{4}{p{0.6\linewidth}}{
        \cellcolor[rgb]{0.9,0.9,0.9}{
            \makecell[{{p{\linewidth}}}]{
                \texttt{\tiny{[P1$\rangle$GM]}}
                \texttt{PREMAKNI SE: L, (2, 1)} \\
            }
        }
    }
    & & \\ \\

    \theutterance \stepcounter{utterance}  
    & & & \multicolumn{2}{p{0.3\linewidth}}{
        \cellcolor[rgb]{0.9,0.9,0.9}{
            \makecell[{{p{\linewidth}}}]{
                \texttt{\tiny{[GM$|$GM]}}
                \texttt{PREMAKNI SE: L, (2, 1)} \\
            }
        }
    }
    & & \\ \\

    \theutterance \stepcounter{utterance}  
    & & & \multicolumn{2}{p{0.3\linewidth}}{
        \cellcolor[rgb]{0.9,0.9,0.9}{
            \makecell[{{p{\linewidth}}}]{
                \texttt{\tiny{[GM$|$GM]}}
                \texttt{Kazen: (2,1) ni prazen, vendar vsebuje '│'.} \\
\\ 
\texttt{Skupaj ste zbrali 7 od 8 kazni.} \\
\texttt{Prosim, poskusite znova!} \\
            }
        }
    }
    & & \\ \\

    \theutterance \stepcounter{utterance}  
    & & \multicolumn{4}{p{0.6\linewidth}}{
        \cellcolor[rgb]{0.9,0.9,0.9}{
            \makecell[{{p{\linewidth}}}]{
                \texttt{\tiny{[P1$\langle$GM]}}
                \texttt{Kazen: (2,1) ni prazen, vendar vsebuje '│'.} \\
\\ 
\texttt{Skupaj ste zbrali 7 od 8 kazni.} \\
\texttt{Prosim, poskusite znova!} \\
            }
        }
    }
    & \\ \\

    \theutterance \stepcounter{utterance}  
    & \multicolumn{4}{p{0.6\linewidth}}{
        \cellcolor[rgb]{0.9,0.9,0.9}{
            \makecell[{{p{\linewidth}}}]{
                \texttt{\tiny{[P1$\rangle$GM]}}
                \texttt{RECI: Oprostite, moja napaka pri branju koordinat in ovir. Moram biti bolj previden s simboli. Ker sem že pri 7 kaznih, ne smem več delati napak. Prosim, bodite potrpežljivi, poskušam najti prazne celice.} \\
            }
        }
    }
    & & \\ \\

    \theutterance \stepcounter{utterance}  
    & & & \multicolumn{2}{p{0.3\linewidth}}{
        \cellcolor[rgb]{0.9,0.9,0.9}{
            \makecell[{{p{\linewidth}}}]{
                \texttt{\tiny{[GM$|$GM]}}
                \texttt{RECI: Oprostite, moja napaka pri branju koordinat in ovir. Moram biti bolj previden s simboli. Ker sem že pri 7 kaznih, ne smem več delati napak. Prosim, bodite potrpežljivi, poskušam najti prazne celice.} \\
            }
        }
    }
    & & \\ \\

    \theutterance \stepcounter{utterance}  
    & & \multicolumn{4}{p{0.6\linewidth}}{
        \cellcolor[rgb]{0.9,0.9,0.9}{
            \makecell[{{p{\linewidth}}}]{
                \texttt{\tiny{[GM$\rangle$P2]}}
                \texttt{Vaše sporočilo je bilo posredovano drugemu igralcu.} \\
\\ 
\\ 
\texttt{Trenutno igrate krog 4 od največ 12.} \\
\texttt{Skupaj ste zbrali 7 od 8 kazni.} \\
\texttt{Drugi igralec je poslal to sporočilo:} \\
\\ 
\\ 
\texttt{"Oprostite, moja napaka pri branju koordinat in ovir. Moram biti bolj previden s simboli. Ker sem že pri 7 kaznih, ne smem več delati napak. Prosim, bodite potrpežljivi, poskušam najti prazne celice."} \\
\\ 
\\ 
\texttt{Kakšen je vaš naslednji ukaz?} \\
            }
        }
    }
    & \\ \\

    \theutterance \stepcounter{utterance}  
    & & & \multicolumn{4}{p{0.6\linewidth}}{
        \cellcolor[rgb]{0.9,0.9,0.9}{
            \makecell[{{p{\linewidth}}}]{
                \texttt{\tiny{[GM$\langle$P2]}}
                \texttt{RECI: Razumem. Ne smemo tvegati več kazni. Jaz bom nadaljeval z premikanjem svojih predmetov proti cilju.} \\
            }
        }
    }
     \\ \\

    \theutterance \stepcounter{utterance}  
    & & & \multicolumn{2}{p{0.3\linewidth}}{
        \cellcolor[rgb]{0.9,0.9,0.9}{
            \makecell[{{p{\linewidth}}}]{
                \texttt{\tiny{[GM$|$GM]}}
                \texttt{RECI: Razumem. Ne smemo tvegati več kazni. Jaz bom nadaljeval z premikanjem svojih predmetov proti cilju.} \\
            }
        }
    }
    & & \\ \\

    \theutterance \stepcounter{utterance}  
    & & \multicolumn{4}{p{0.6\linewidth}}{
        \cellcolor[rgb]{0.9,0.9,0.9}{
            \makecell[{{p{\linewidth}}}]{
                \texttt{\tiny{[P1$\langle$GM]}}
                \texttt{Vaše sporočilo je bilo posredovano drugemu igralcu.} \\
\\ 
\\ 
\texttt{Trenutno igrate krog 4 od največ 12.} \\
\texttt{Skupaj ste zbrali 7 od 8 kazni.} \\
\texttt{Drugi igralec je poslal to sporočilo:} \\
\\ 
\\ 
\texttt{"Razumem. Ne smemo tvegati več kazni. Jaz bom nadaljeval z premikanjem svojih predmetov proti cilju."} \\
\\ 
\\ 
\texttt{Kakšen je vaš naslednji ukaz?} \\
            }
        }
    }
    & \\ \\

    \theutterance \stepcounter{utterance}  
    & \multicolumn{4}{p{0.6\linewidth}}{
        \cellcolor[rgb]{0.9,0.9,0.9}{
            \makecell[{{p{\linewidth}}}]{
                \texttt{\tiny{[P1$\rangle$GM]}}
                \texttt{RECI: Hvala za razumevanje. Ker imam še samo eno možnost brez izgube igre, ne bom več poskušal premikati predmetov, dokler ne bo popolnoma jasno, katera celica je zares prazna (brez çizk). Prosim, nadaljujte vi, jaz pa bom samo spremljal vaš napredek.} \\
            }
        }
    }
    & & \\ \\

    \theutterance \stepcounter{utterance}  
    & & & \multicolumn{2}{p{0.3\linewidth}}{
        \cellcolor[rgb]{0.9,0.9,0.9}{
            \makecell[{{p{\linewidth}}}]{
                \texttt{\tiny{[GM$|$GM]}}
                \texttt{RECI: Hvala za razumevanje. Ker imam še samo eno možnost brez izgube igre, ne bom več poskušal premikati predmetov, dokler ne bo popolnoma jasno, katera celica je zares prazna (brez çizk). Prosim, nadaljujte vi, jaz pa bom samo spremljal vaš napredek.} \\
            }
        }
    }
    & & \\ \\

    \theutterance \stepcounter{utterance}  
    & & \multicolumn{4}{p{0.6\linewidth}}{
        \cellcolor[rgb]{0.9,0.9,0.9}{
            \makecell[{{p{\linewidth}}}]{
                \texttt{\tiny{[GM$\rangle$P2]}}
                \texttt{Vaše sporočilo je bilo posredovano drugemu igralcu.} \\
\\ 
\\ 
\texttt{Trenutno igrate krog 5 od največ 12.} \\
\texttt{Skupaj ste zbrali 7 od 8 kazni.} \\
\texttt{Drugi igralec je poslal to sporočilo:} \\
\\ 
\\ 
\texttt{"Hvala za razumevanje. Ker imam še samo eno možnost brez izgube igre, ne bom več poskušal premikati predmetov, dokler ne bo popolnoma jasno, katera celica je zares prazna (brez çizk). Prosim, nadaljujte vi, jaz pa bom samo spremljal vaš napredek."} \\
\\ 
\\ 
\texttt{Kakšen je vaš naslednji ukaz?} \\
            }
        }
    }
    & \\ \\

    \theutterance \stepcounter{utterance}  
    & & & \multicolumn{4}{p{0.6\linewidth}}{
        \cellcolor[rgb]{0.9,0.9,0.9}{
            \makecell[{{p{\linewidth}}}]{
                \texttt{\tiny{[GM$\langle$P2]}}
                \texttt{PREMAKNI SE: P, (3, 1)} \\
            }
        }
    }
     \\ \\

    \theutterance \stepcounter{utterance}  
    & & & \multicolumn{2}{p{0.3\linewidth}}{
        \cellcolor[rgb]{0.9,0.9,0.9}{
            \makecell[{{p{\linewidth}}}]{
                \texttt{\tiny{[GM$|$GM]}}
                \texttt{PREMAKNI SE: P, (3, 1)} \\
            }
        }
    }
    & & \\ \\

    \theutterance \stepcounter{utterance}  
    & & & \multicolumn{2}{p{0.3\linewidth}}{
        \cellcolor[rgb]{0.9,0.9,0.9}{
            \makecell[{{p{\linewidth}}}]{
                \texttt{\tiny{[GM$|$GM]}}
                \texttt{P moved to (3, 1)} \\
            }
        }
    }
    & & \\ \\

    \theutterance \stepcounter{utterance}  
    & & \multicolumn{4}{p{0.6\linewidth}}{
        \cellcolor[rgb]{0.9,0.9,0.9}{
            \makecell[{{p{\linewidth}}}]{
                \texttt{\tiny{[P1$\langle$GM]}}
                \texttt{Vaše sporočilo je bilo posredovano drugemu igralcu.} \\
\\ 
\\ 
\texttt{Trenutno igrate krog 5 od največ 12.} \\
\texttt{Skupaj ste zbrali 7 od 8 kazni.} \\
\texttt{Drugi igralec je premaknil predmet na svoji mreži.} \\
\\ 
\\ 
\texttt{Kakšen je vaš naslednji ukaz?} \\
            }
        }
    }
    & \\ \\

    \theutterance \stepcounter{utterance}  
    & \multicolumn{4}{p{0.6\linewidth}}{
        \cellcolor[rgb]{0.9,0.9,0.9}{
            \makecell[{{p{\linewidth}}}]{
                \texttt{\tiny{[P1$\rangle$GM]}}
                \texttt{RECI: Kakšen je vaš trenutni položaj predmetov? Želim vedeti, kje so vaši L, C in P, da bom lahko v naslednjem krogu (če bo mogoče) poskušal ugotoviti varne koordinate na svoji mreži.} \\
            }
        }
    }
    & & \\ \\

    \theutterance \stepcounter{utterance}  
    & & & \multicolumn{2}{p{0.3\linewidth}}{
        \cellcolor[rgb]{0.9,0.9,0.9}{
            \makecell[{{p{\linewidth}}}]{
                \texttt{\tiny{[GM$|$GM]}}
                \texttt{RECI: Kakšen je vaš trenutni položaj predmetov? Želim vedeti, kje so vaši L, C in P, da bom lahko v naslednjem krogu (če bo mogoče) poskušal ugotoviti varne koordinate na svoji mreži.} \\
            }
        }
    }
    & & \\ \\

    \theutterance \stepcounter{utterance}  
    & & \multicolumn{4}{p{0.6\linewidth}}{
        \cellcolor[rgb]{0.9,0.9,0.9}{
            \makecell[{{p{\linewidth}}}]{
                \texttt{\tiny{[GM$\rangle$P2]}}
                \texttt{Uspešno premaknjen 'P' v (3,1). Vaša posodobljena mreža je videti takole:} \\
\texttt{```} \\
\texttt{ 1234567} \\
\texttt{╔═╤═╤═══╗} \\
\texttt{║L│P│◌◌◌║ 1} \\
\texttt{║◌├─┤◌◌◌║ 2} \\
\texttt{╟─┤◌├─┬─╢ 3} \\
\texttt{║◌│◌│◌│◌║ 4} \\
\texttt{╟─┼─┘C│◌║ 5} \\
\texttt{║◌│◌◌┌┴─╢ 6} \\
\texttt{║◌│◌◌│◌◌║ 7} \\
\texttt{╚═╧══╧══╝} \\
\texttt{```} \\
\texttt{Trenutno igrate krog 6 od največ 12.} \\
\texttt{Skupaj ste zbrali 7 od 8 kazni.} \\
\texttt{Drugi igralec je poslal to sporočilo:} \\
\\ 
\\ 
\texttt{"Kakšen je vaš trenutni položaj predmetov? Želim vedeti, kje so vaši L, C in P, da bom lahko v naslednjem krogu (če bo mogoče) poskušal ugotoviti varne koordinate na svoji mreži."} \\
\\ 
\\ 
\texttt{Kakšen je vaš naslednji ukaz?} \\
            }
        }
    }
    & \\ \\

    \theutterance \stepcounter{utterance}  
    & & & \multicolumn{4}{p{0.6\linewidth}}{
        \cellcolor[rgb]{0.9,0.9,0.9}{
            \makecell[{{p{\linewidth}}}]{
                \texttt{\tiny{[GM$\langle$P2]}}
                \texttt{RECI: Moji predmeti so trenutno: L na (1,1), P na (3,1) in C na (4,5).} \\
            }
        }
    }
     \\ \\

    \theutterance \stepcounter{utterance}  
    & & & \multicolumn{2}{p{0.3\linewidth}}{
        \cellcolor[rgb]{0.9,0.9,0.9}{
            \makecell[{{p{\linewidth}}}]{
                \texttt{\tiny{[GM$|$GM]}}
                \texttt{RECI: Moji predmeti so trenutno: L na (1,1), P na (3,1) in C na (4,5).} \\
            }
        }
    }
    & & \\ \\

    \theutterance \stepcounter{utterance}  
    & & \multicolumn{4}{p{0.6\linewidth}}{
        \cellcolor[rgb]{0.9,0.9,0.9}{
            \makecell[{{p{\linewidth}}}]{
                \texttt{\tiny{[P1$\langle$GM]}}
                \texttt{Vaše sporočilo je bilo posredovano drugemu igralcu.} \\
\\ 
\\ 
\texttt{Trenutno igrate krog 6 od največ 12.} \\
\texttt{Skupaj ste zbrali 7 od 8 kazni.} \\
\texttt{Drugi igralec je poslal to sporočilo:} \\
\\ 
\\ 
\texttt{"Moji predmeti so trenutno: L na (1,1), P na (3,1) in C na (4,5)."} \\
\\ 
\\ 
\texttt{Kakšen je vaš naslednji ukaz?} \\
            }
        }
    }
    & \\ \\

    \theutterance \stepcounter{utterance}  
    & \multicolumn{4}{p{0.6\linewidth}}{
        \cellcolor[rgb]{0.9,0.9,0.9}{
            \makecell[{{p{\linewidth}}}]{
                \texttt{\tiny{[P1$\rangle$GM]}}
                \texttt{RECI: Moja trenutna situacija je: L na (1,1), P na (3,1) in C na (1,2). Ker sem pri 7 kaznih, ne bom več izvajal ukazov PREMAKNI SE, razen če bom popolnoma prepričan o koordinatah. Ker so vaši L in P že na ciljnih mestih (1,1) in (3,1), se bomo lahko ujemali.} \\
            }
        }
    }
    & & \\ \\

    \theutterance \stepcounter{utterance}  
    & & & \multicolumn{2}{p{0.3\linewidth}}{
        \cellcolor[rgb]{0.9,0.9,0.9}{
            \makecell[{{p{\linewidth}}}]{
                \texttt{\tiny{[GM$|$GM]}}
                \texttt{RECI: Moja trenutna situacija je: L na (1,1), P na (3,1) in C na (1,2). Ker sem pri 7 kaznih, ne bom več izvajal ukazov PREMAKNI SE, razen če bom popolnoma prepričan o koordinatah. Ker so vaši L in P že na ciljnih mestih (1,1) in (3,1), se bomo lahko ujemali.} \\
            }
        }
    }
    & & \\ \\

    \theutterance \stepcounter{utterance}  
    & & \multicolumn{4}{p{0.6\linewidth}}{
        \cellcolor[rgb]{0.9,0.9,0.9}{
            \makecell[{{p{\linewidth}}}]{
                \texttt{\tiny{[GM$\rangle$P2]}}
                \texttt{Vaše sporočilo je bilo posredovano drugemu igralcu.} \\
\\ 
\\ 
\texttt{Trenutno igrate krog 7 od največ 12.} \\
\texttt{Skupaj ste zbrali 7 od 8 kazni.} \\
\texttt{Drugi igralec je poslal to sporočilo:} \\
\\ 
\\ 
\texttt{"Moja trenutna situacija je: L na (1,1), P na (3,1) in C na (1,2). Ker sem pri 7 kaznih, ne bom več izvajal ukazov PREMAKNI SE, razen če bom popolnoma prepričan o koordinatah. Ker so vaši L in P že na ciljnih mestih (1,1) in (3,1), se bomo lahko ujemali."} \\
\\ 
\\ 
\texttt{Kakšen je vaš naslednji ukaz?} \\
            }
        }
    }
    & \\ \\

    \theutterance \stepcounter{utterance}  
    & & & \multicolumn{4}{p{0.6\linewidth}}{
        \cellcolor[rgb]{0.9,0.9,0.9}{
            \makecell[{{p{\linewidth}}}]{
                \texttt{\tiny{[GM$\langle$P2]}}
                \texttt{PREMAKNI SE: C, (2, 1)} \\
            }
        }
    }
     \\ \\

    \theutterance \stepcounter{utterance}  
    & & & \multicolumn{2}{p{0.3\linewidth}}{
        \cellcolor[rgb]{0.9,0.9,0.9}{
            \makecell[{{p{\linewidth}}}]{
                \texttt{\tiny{[GM$|$GM]}}
                \texttt{PREMAKNI SE: C, (2, 1)} \\
            }
        }
    }
    & & \\ \\

    \theutterance \stepcounter{utterance}  
    & & & \multicolumn{2}{p{0.3\linewidth}}{
        \cellcolor[rgb]{0.9,0.9,0.9}{
            \makecell[{{p{\linewidth}}}]{
                \texttt{\tiny{[GM$|$GM]}}
                \texttt{Kazen: (2,1) ni prazen, vendar vsebuje '│'.} \\
\\ 
\texttt{Skupaj ste zbrali 8 od 8 kazni.} \\
\texttt{Prosim, poskusite znova!} \\
            }
        }
    }
    & & \\ \\

    \theutterance \stepcounter{utterance}  
    & & \multicolumn{4}{p{0.6\linewidth}}{
        \cellcolor[rgb]{0.9,0.9,0.9}{
            \makecell[{{p{\linewidth}}}]{
                \texttt{\tiny{[GM$\rangle$P2]}}
                \texttt{Kazen: (2,1) ni prazen, vendar vsebuje '│'.} \\
\\ 
\texttt{Skupaj ste zbrali 8 od 8 kazni.} \\
\texttt{Prosim, poskusite znova!} \\
            }
        }
    }
    & \\ \\

    \theutterance \stepcounter{utterance}  
    & & & \multicolumn{4}{p{0.6\linewidth}}{
        \cellcolor[rgb]{0.9,0.9,0.9}{
            \makecell[{{p{\linewidth}}}]{
                \texttt{\tiny{[GM$\langle$P2]}}
                \texttt{PREMAKNI SE: C, (2, 2)} \\
            }
        }
    }
     \\ \\

    \theutterance \stepcounter{utterance}  
    & & & \multicolumn{2}{p{0.3\linewidth}}{
        \cellcolor[rgb]{0.9,0.9,0.9}{
            \makecell[{{p{\linewidth}}}]{
                \texttt{\tiny{[GM$|$GM]}}
                \texttt{PREMAKNI SE: C, (2, 2)} \\
            }
        }
    }
    & & \\ \\

    \theutterance \stepcounter{utterance}  
    & & & \multicolumn{2}{p{0.3\linewidth}}{
        \cellcolor[rgb]{0.9,0.9,0.9}{
            \makecell[{{p{\linewidth}}}]{
                \texttt{\tiny{[GM$|$GM]}}
                \texttt{Kazen: (2,2) ni prazen, vendar vsebuje '├'.} \\
\\ 
\texttt{Skupaj ste zbrali 9 od 8 kazni.} \\
\texttt{Prosim, poskusite znova!} \\
            }
        }
    }
    & & \\ \\

    \theutterance \stepcounter{utterance}  
    & & \multicolumn{4}{p{0.6\linewidth}}{
        \cellcolor[rgb]{0.9,0.9,0.9}{
            \makecell[{{p{\linewidth}}}]{
                \texttt{\tiny{[GM$\rangle$P2]}}
                \texttt{Kazen: (2,2) ni prazen, vendar vsebuje '├'.} \\
\\ 
\texttt{Skupaj ste zbrali 9 od 8 kazni.} \\
\texttt{Prosim, poskusite znova!} \\
            }
        }
    }
    & \\ \\

    \theutterance \stepcounter{utterance}  
    & & & \multicolumn{4}{p{0.6\linewidth}}{
        \cellcolor[rgb]{0.9,0.9,0.9}{
            \makecell[{{p{\linewidth}}}]{
                \texttt{\tiny{[GM$\langle$P2]}}
                \texttt{PREMAKNI SE: C, (4, 4)} \\
            }
        }
    }
     \\ \\

    \theutterance \stepcounter{utterance}  
    & & & \multicolumn{2}{p{0.3\linewidth}}{
        \cellcolor[rgb]{0.9,0.9,0.9}{
            \makecell[{{p{\linewidth}}}]{
                \texttt{\tiny{[GM$|$GM]}}
                \texttt{PREMAKNI SE: C, (4, 4)} \\
            }
        }
    }
    & & \\ \\

    \theutterance \stepcounter{utterance}  
    & & & \multicolumn{2}{p{0.3\linewidth}}{
        \cellcolor[rgb]{0.9,0.9,0.9}{
            \makecell[{{p{\linewidth}}}]{
                \texttt{\tiny{[GM$|$GM]}}
                \texttt{Kazen: (4,4) ni prazen, vendar vsebuje '│'.} \\
\\ 
\texttt{Skupaj ste zbrali 10 od 8 kazni.} \\
\texttt{Prosim, poskusite znova!} \\
            }
        }
    }
    & & \\ \\

    \theutterance \stepcounter{utterance}  
    & & \multicolumn{4}{p{0.6\linewidth}}{
        \cellcolor[rgb]{0.9,0.9,0.9}{
            \makecell[{{p{\linewidth}}}]{
                \texttt{\tiny{[GM$\rangle$P2]}}
                \texttt{Kazen: (4,4) ni prazen, vendar vsebuje '│'.} \\
\\ 
\texttt{Skupaj ste zbrali 10 od 8 kazni.} \\
\texttt{Prosim, poskusite znova!} \\
            }
        }
    }
    & \\ \\

    \theutterance \stepcounter{utterance}  
    & & & \multicolumn{4}{p{0.6\linewidth}}{
        \cellcolor[rgb]{0.9,0.9,0.9}{
            \makecell[{{p{\linewidth}}}]{
                \texttt{\tiny{[GM$\langle$P2]}}
                \texttt{PREMAKNI SE: C, (5, 5)} \\
            }
        }
    }
     \\ \\

    \theutterance \stepcounter{utterance}  
    & & & \multicolumn{2}{p{0.3\linewidth}}{
        \cellcolor[rgb]{0.9,0.9,0.9}{
            \makecell[{{p{\linewidth}}}]{
                \texttt{\tiny{[GM$|$GM]}}
                \texttt{PREMAKNI SE: C, (5, 5)} \\
            }
        }
    }
    & & \\ \\

    \theutterance \stepcounter{utterance}  
    & & & \multicolumn{2}{p{0.3\linewidth}}{
        \cellcolor[rgb]{0.9,0.9,0.9}{
            \makecell[{{p{\linewidth}}}]{
                \texttt{\tiny{[GM$|$GM]}}
                \texttt{Kazen: (5,5) ni prazen, vendar vsebuje 'C'.} \\
\\ 
\texttt{Skupaj ste zbrali 11 od 8 kazni.} \\
\texttt{Prosim, poskusite znova!} \\
            }
        }
    }
    & & \\ \\

    \theutterance \stepcounter{utterance}  
    & & \multicolumn{4}{p{0.6\linewidth}}{
        \cellcolor[rgb]{0.9,0.9,0.9}{
            \makecell[{{p{\linewidth}}}]{
                \texttt{\tiny{[GM$\rangle$P2]}}
                \texttt{Kazen: (5,5) ni prazen, vendar vsebuje 'C'.} \\
\\ 
\texttt{Skupaj ste zbrali 11 od 8 kazni.} \\
\texttt{Prosim, poskusite znova!} \\
            }
        }
    }
    & \\ \\

    \theutterance \stepcounter{utterance}  
    & & & \multicolumn{4}{p{0.6\linewidth}}{
        \cellcolor[rgb]{0.9,0.9,0.9}{
            \makecell[{{p{\linewidth}}}]{
                \texttt{\tiny{[GM$\langle$P2]}}
                \texttt{PREMAKNI SE: C, (6, 5)} \\
            }
        }
    }
     \\ \\

    \theutterance \stepcounter{utterance}  
    & & & \multicolumn{2}{p{0.3\linewidth}}{
        \cellcolor[rgb]{0.9,0.9,0.9}{
            \makecell[{{p{\linewidth}}}]{
                \texttt{\tiny{[GM$|$GM]}}
                \texttt{PREMAKNI SE: C, (6, 5)} \\
            }
        }
    }
    & & \\ \\

    \theutterance \stepcounter{utterance}  
    & & & \multicolumn{2}{p{0.3\linewidth}}{
        \cellcolor[rgb]{0.9,0.9,0.9}{
            \makecell[{{p{\linewidth}}}]{
                \texttt{\tiny{[GM$|$GM]}}
                \texttt{Kazen: (6,5) ni prazen, vendar vsebuje '│'.} \\
\\ 
\texttt{Skupaj ste zbrali 12 od 8 kazni.} \\
\texttt{Prosim, poskusite znova!} \\
            }
        }
    }
    & & \\ \\

    \theutterance \stepcounter{utterance}  
    & & \multicolumn{4}{p{0.6\linewidth}}{
        \cellcolor[rgb]{0.9,0.9,0.9}{
            \makecell[{{p{\linewidth}}}]{
                \texttt{\tiny{[GM$\rangle$P2]}}
                \texttt{Kazen: (6,5) ni prazen, vendar vsebuje '│'.} \\
\\ 
\texttt{Skupaj ste zbrali 12 od 8 kazni.} \\
\texttt{Prosim, poskusite znova!} \\
            }
        }
    }
    & \\ \\

    \theutterance \stepcounter{utterance}  
    & & & \multicolumn{4}{p{0.6\linewidth}}{
        \cellcolor[rgb]{0.9,0.9,0.9}{
            \makecell[{{p{\linewidth}}}]{
                \texttt{\tiny{[GM$\langle$P2]}}
                \texttt{PREMAKNI SE: C, (7, 5)} \\
            }
        }
    }
     \\ \\

    \theutterance \stepcounter{utterance}  
    & & & \multicolumn{2}{p{0.3\linewidth}}{
        \cellcolor[rgb]{0.9,0.9,0.9}{
            \makecell[{{p{\linewidth}}}]{
                \texttt{\tiny{[GM$|$GM]}}
                \texttt{PREMAKNI SE: C, (7, 5)} \\
            }
        }
    }
    & & \\ \\

    \theutterance \stepcounter{utterance}  
    & & & \multicolumn{2}{p{0.3\linewidth}}{
        \cellcolor[rgb]{0.9,0.9,0.9}{
            \makecell[{{p{\linewidth}}}]{
                \texttt{\tiny{[GM$|$GM]}}
                \texttt{C moved to (7, 5)} \\
            }
        }
    }
    & & \\ \\

    \theutterance \stepcounter{utterance}  
    & & \multicolumn{4}{p{0.6\linewidth}}{
        \cellcolor[rgb]{0.9,0.9,0.9}{
            \makecell[{{p{\linewidth}}}]{
                \texttt{\tiny{[P1$\langle$GM]}}
                \texttt{Vaše sporočilo je bilo posredovano drugemu igralcu.} \\
\\ 
\\ 
\texttt{Trenutno igrate krog 7 od največ 12.} \\
\texttt{Skupaj ste zbrali 12 od 8 kazni.} \\
\texttt{Drugi igralec je premaknil predmet na svoji mreži.} \\
\\ 
\\ 
\texttt{Kakšen je vaš naslednji ukaz?} \\
            }
        }
    }
    & \\ \\

    \theutterance \stepcounter{utterance}  
    & \multicolumn{4}{p{0.6\linewidth}}{
        \cellcolor[rgb]{0.9,0.9,0.9}{
            \makecell[{{p{\linewidth}}}]{
                \texttt{\tiny{[P1$\rangle$GM]}}
                \texttt{Igra je končana, ker smo presegli omejitev 8 kazni. Moja zadnja strategija je bila, da ne tavam več z premiki, da ne bi povzročil izgube, vendar je bila igra že izgubljena zaradi predhodnih napak pri branju mreže.} \\
            }
        }
    }
    & & \\ \\

    \theutterance \stepcounter{utterance}  
    & & & \multicolumn{2}{p{0.3\linewidth}}{
        \cellcolor[rgb]{0.9,0.9,0.9}{
            \makecell[{{p{\linewidth}}}]{
                \texttt{\tiny{[GM$|$GM]}}
                \texttt{Igra je končana, ker smo presegli omejitev 8 kazni. Moja zadnja strategija je bila, da ne tavam več z premiki, da ne bi povzročil izgube, vendar je bila igra že izgubljena zaradi predhodnih napak pri branju mreže.} \\
            }
        }
    }
    & & \\ \\

    \theutterance \stepcounter{utterance}  
    & & & \multicolumn{2}{p{0.3\linewidth}}{
        \cellcolor[rgb]{0.9,0.9,0.9}{
            \makecell[{{p{\linewidth}}}]{
                \texttt{\tiny{[GM$|$GM]}}
                \texttt{Invalid response format} \\
            }
        }
    }
    & & \\ \\

    \theutterance \stepcounter{utterance}  
    & & & \multicolumn{2}{p{0.3\linewidth}}{
        \cellcolor[rgb]{0.9,0.9,0.9}{
            \makecell[{{p{\linewidth}}}]{
                \texttt{\tiny{[GM$|$GM]}}
                \texttt{Maximum number of penalties exceeded} \\
            }
        }
    }
    & & \\ \\

    \theutterance \stepcounter{utterance}  
    & & & \multicolumn{2}{p{0.3\linewidth}}{
        \cellcolor[rgb]{0.9,0.9,0.9}{
            \makecell[{{p{\linewidth}}}]{
                \texttt{\tiny{[GM$|$GM]}}
                \texttt{Initial states:} \\
\texttt{```} \\
\texttt{Player 1:} \\
\texttt{ 1234567} \\
\texttt{╔═╤═╤═══╗} \\
\texttt{║L│◌│◌◌◌║ 1} \\
\texttt{║C├─┤◌◌◌║ 2} \\
\texttt{╟─┤◌├─┬─╢ 3} \\
\texttt{║◌│◌│◌│◌║ 4} \\
\texttt{╟─┼─┘◌│◌║ 5} \\
\texttt{║◌│◌◌┌┴─╢ 6} \\
\texttt{║◌│◌◌│◌P║ 7} \\
\texttt{╚═╧══╧══╝} \\
\\ 
\\ 
\texttt{Player 2:} \\
\texttt{ 1234567} \\
\texttt{╔═╤═╤═══╗} \\
\texttt{║◌│◌│◌◌◌║ 1} \\
\texttt{║◌├─┤◌L◌║ 2} \\
\texttt{╟─┤P├─┬─╢ 3} \\
\texttt{║◌│◌│◌│◌║ 4} \\
\texttt{╟─┼─┘C│◌║ 5} \\
\texttt{║◌│◌◌┌┴─╢ 6} \\
\texttt{║◌│◌◌│◌◌║ 7} \\
\texttt{╚═╧══╧══╝} \\
\\ 
\texttt{```} \\
            }
        }
    }
    & & \\ \\

    \theutterance \stepcounter{utterance}  
    & & & \multicolumn{2}{p{0.3\linewidth}}{
        \cellcolor[rgb]{0.9,0.9,0.9}{
            \makecell[{{p{\linewidth}}}]{
                \texttt{\tiny{[GM$|$GM]}}
                \texttt{End states:} \\
\texttt{```} \\
\texttt{Player 1:} \\
\texttt{ 1234567} \\
\texttt{╔═╤═╤═══╗} \\
\texttt{║L│P│◌◌◌║ 1} \\
\texttt{║C├─┤◌◌◌║ 2} \\
\texttt{╟─┤◌├─┬─╢ 3} \\
\texttt{║◌│◌│◌│◌║ 4} \\
\texttt{╟─┼─┘◌│◌║ 5} \\
\texttt{║◌│◌◌┌┴─╢ 6} \\
\texttt{║◌│◌◌│◌◌║ 7} \\
\texttt{╚═╧══╧══╝} \\
\\ 
\\ 
\texttt{Player 2:} \\
\texttt{ 1234567} \\
\texttt{╔═╤═╤═══╗} \\
\texttt{║L│P│◌◌◌║ 1} \\
\texttt{║◌├─┤◌◌◌║ 2} \\
\texttt{╟─┤◌├─┬─╢ 3} \\
\texttt{║◌│◌│◌│◌║ 4} \\
\texttt{╟─┼─┘◌│C║ 5} \\
\texttt{║◌│◌◌┌┴─╢ 6} \\
\texttt{║◌│◌◌│◌◌║ 7} \\
\texttt{╚═╧══╧══╝} \\
\\ 
\texttt{```} \\
            }
        }
    }
    & & \\ \\

    \theutterance \stepcounter{utterance}  
    & & & \multicolumn{2}{p{0.3\linewidth}}{
        \cellcolor[rgb]{0.9,0.9,0.9}{
            \makecell[{{p{\linewidth}}}]{
                \texttt{\tiny{[GM$|$GM]}}
                \texttt{* success: False} \\
\texttt{* lose: False} \\
\texttt{* aborted: True} \\
\texttt{{-}{-}{-}{-}{-}{-}{-}} \\
\texttt{* Shifts: 3.00} \\
\texttt{* Max Shifts: 4.00} \\
\texttt{* Min Shifts: 2.00} \\
\texttt{* End Distance Sum: 6.71} \\
\texttt{* Init Distance Sum: 15.76} \\
\texttt{* Expected Distance Sum: 12.57} \\
\texttt{* Penalties: 13.00} \\
\texttt{* Max Penalties: 8.00} \\
\texttt{* Rounds: 7.00} \\
\texttt{* Max Rounds: 12.00} \\
\texttt{* Object Count: 3.00} \\
            }
        }
    }
    & & \\ \\

\end{supertabular}
}

\end{document}
